\section{Research Approach}
\label{sec:research-approach}

This thesis is a platform engineering thesis: its primary contribution is a designed and implemented software artifact (a researcher-operated adaptive reading platform) rather than an empirical study of human reading behaviour.
The research approach is therefore shaped by the nature of that contribution: the goal is not to measure a natural phenomenon but to construct, justify, and evaluate a designed system.
The Double Diamond model \cite{designcouncil2023diamond}, introduced at the opening of this chapter, serves as the overarching process framework, organising the work into two sequential diverge-then-converge cycles that cover the problem space and the solution space respectively.

This artifact orientation determines what counts as a research output in this thesis.
Design decisions are first-class contributions that must be motivated, documented, and evaluated, not merely implemented.
The architecture itself (the separation of sensing, decision-making, intervention execution, and researcher interface into independently replaceable modules) is the central claim this thesis makes, and the evaluation is designed to assess whether that claim holds in the implemented system.

The research questions stated in \cref{sec:final-problem-statement} reflect this orientation directly.
RQ1 through RQ4 all ask about properties of the designed system (modularity, sensing pipeline feasibility, context-preserving intervention, and researcher experimentability), not properties of human reading behaviour.
This is a deliberate boundary: whether specific typographic interventions measurably improve reading outcomes for participants is a question for a future empirical study conducted with this platform, and it is explicitly outside the scope of this work.
The platform is the instrument that makes such future studies possible; building and validating that instrument is the contribution of this thesis.
